\section{ПОСТАНОВКА ЗАДАЧИ (INVOICES-S13)}

\subsection{Назначение информационной системы}
Информационная система \texttt{invoices-s13} разработана для автоматизации учета финансовых документов (счетов) в рамках микросервисной архитектуры. Основная цель системы --- обеспечить надежное хранение данных о счетах и предоставить гибкий интерфейс для их создания и получения внешними потребителями.

\subsection{Функциональные требования к системе}
Система должна реализовывать следующий набор функций:
\begin{itemize}
    \item Создание записи о счете с обязательным указанием суммы (\texttt{amount}).
    \item Получение подробной информации по идентификатору счета.
    \item Формирование списка всех существующих счетов через потоковую передачу данных.
    \item Корректная трансляция данных между различными форматами представления (Protobuf и JSON).
\end{itemize}

\subsection{Нефункциональные требования}
\begin{itemize}
    \item \textbf{Стек технологий}: Использование Python для всех компонентов системы.
    \item \textbf{Взаимодействие}: Применение gRPC для внутренних вызовов и REST для внешних.
    \item \textbf{Изоляция}: Каждый компонент должен функционировать в собственном Docker-контейнере.
\end{itemize}

\subsection{Спецификация варианта s13}
Согласно заданию, система работает с объектом \texttt{invoice}. Основной числовой атрибут --- \texttt{amount} (тип \texttt{float}). Все компоненты системы используют идентификатор \texttt{s13} в названиях и конфигурациях.
